\section*{Capítulo \ref{ch:b1:predominant-reaction} --- Equilíbrios ácido--base e o método da reação predominante}

\begin{solution}{exo:b1:predominant-reaction:1}
Bases conjugadas: \ce{HSO4-}, \ce{CO3^{2-}}, \ce{NH3}, \ce{OH-}. Ácidos
conjugados: \ce{NH4+}, \ce{H2PO4-}, \ce{H2O}, \ce{H3O+}. \omterm{def:b1:predominant-reaction:ampholyte}{Anfólitos}:
\ce{HCO3-}, \ce{H2O}, \ce{HPO4^{2-}} (e \ce{HSO4-}, base do ácido
sulfúrico e ácido de seu segundo par).
\end{solution}

\begin{solution}{exo:b1:predominant-reaction:2}
\ce{HF} ($K_a = \num{6.9e-4}$) $>$ \ce{CH3COOH} (\num{1.7e-5}) $>$ \ce{HClO}
(\num{2.8e-8}) $>$ \ce{NH4+} (\num{5.6e-10}).
\end{solution}

\begin{solution}{exo:b1:predominant-reaction:3}
\omterm{def:b1:predominant-reaction:strong-acid}{Ácido forte}: $\mathrm{pH} = 2.00$. \omterm{def:b1:predominant-reaction:strong-acid}{Base forte}: $[\ce{OH-}] = 10^{-2}$,
$\mathrm{pH} = 14.00 - 2.00 = 12.00$.
\end{solution}

\begin{solution}{exo:b1:predominant-reaction:4}
\ce{CH3COOH} abaixo de pH 4.76, \ce{CH3COO-} acima. Em pH 3, a forma
ácida predomina (98\,\%); no sangue (pH 7.4), o íon etanoato. Em pH 3.76,
a fração ácida é $1/(1 + 10^{-1}) = 0.91$.
\end{solution}

\begin{solution}{exo:b1:predominant-reaction:5}
\omterm{def:b1:predominant-reaction:predominant-reaction}{Reação predominante} \ce{NH3 + H2O <=> NH4+ + OH-}, $K = 10^{-(14.00 -
9.25)} = 10^{-4.75}$. Se pouco \ce{NH3} reage: $[\ce{OH-}] = \sqrt{K C} =
\qty{1.3e-3}{mol/L}$, $\mathrm{pH} = 11.13$. Verificação: $[\ce{NH4+}] = 1.3\,\%$
de $C$, e $\mathrm{pH} > \mathrm{p}K_a + 1$: válido.
\end{solution}

\begin{solution}{exo:b1:predominant-reaction:6}
$\mathrm{pH} = \frac12(4.76 + 2.00) = 3.38$ (o valor exato é 3.39);
$\alpha = 10^{-3.38}/0.010 = 0.042$: 4\,\% dissociado, abaixo de 10\,\%,
e $\mathrm{pH} < \mathrm{p}K_a - 1$: válido.
\end{solution}

\begin{solution}{exo:b1:predominant-reaction:7}
\ce{HF + OH- -> F- + H2O}, $K = 10^{14.00 - 3.16} = 10^{10.84}$: total. A
\omterm{def:b1:predominant-reaction:predominant-reaction}{solução equivalente} contém \qty{0.050}{mol/L} de \ce{HF} e de \ce{F-}, um
tampão: $\mathrm{pH} = \mathrm{p}K_a = 3.16$.
\end{solution}

\begin{solution}{exo:b1:predominant-reaction:8}
$\mathrm{pH} = 4.76$. Com o ácido clorídrico, \qty{0.010}{mol} de
etanoato passam a ácido: $\mathrm{pH} = 4.76 + \log(0.040/0.060) = 4.58$, uma
variação de 0.18. Em água pura, o pH cairia de 7.00 para 2.00.
\end{solution}

\begin{solution}{exo:b1:predominant-reaction:9}
Na água, os dois ácidos reagem totalmente, dando \ce{H3O+}: o solvente
os nivela. O ácido etanoico é uma base muito mais fraca que a água, de
modo que nele os dois ácidos só reagem parcialmente, em graus
diferentes, e suas forças podem ser comparadas. O ácido mais forte que
pode existir na água é \ce{H3O+}.
\end{solution}

\begin{solution}{exo:b1:predominant-reaction:10}
\ce{CO3^{2-} + H2O <=> HCO3- + OH-}, $K = 10^{-(14.00 - 10.33)} = 10^{-3.67}$;
$[\ce{OH-}] = \sqrt{10^{-3.67} \times 0.10} = \qty{4.6e-3}{mol/L}$ (4.6\,\%
de $C$), $\mathrm{pH} = 11.67$. A segunda basicidade, \ce{HCO3- + H2O <=>
CO2 + H2O + OH-}, tem $K = 10^{-7.63}$: ela dá $[\ce{CO2}] \approx
10^{-7.6}$, desprezível.
\end{solution}

\begin{solution}{exo:b1:predominant-reaction:11}
$[\mathrm{A}] = [\ce{H3O+}] = \alpha C$ e $[\mathrm{HA}] = (1 - \alpha)C$,
logo $K_a = \alpha^2C/(1 - \alpha)$. Com $C = 10^{-3}$: $\alpha^2 +
0.0174\,\alpha - 0.0174 = 0$, $\alpha = 0.12$. O ácido está 12\,\%
dissociado: a hipótese de um ácido pouco dissociado ($\alpha <
10\,\%$) falha, e é preciso resolver a equação do segundo grau.
\end{solution}

\begin{solution}{exo:b1:predominant-reaction:12}
O ácido clorídrico fixa $[\ce{H3O+}] = \qty{0.010}{mol/L}$: pH 2.00. Então
\[
  [\ce{CH3COO-}] = \frac{K_a[\ce{CH3COOH}]}{[\ce{H3O+}]} = \frac{1.74 \times 10^{-5}
  \times 0.10}{0.010} = \qty{1.7e-4}{mol/L},
\]
desprezível diante de 0.010: o
\ce{H3O+} do \omterm{def:b1:predominant-reaction:strong-acid}{ácido forte} faz recuar o equilíbrio do \omterm{def:b1:predominant-reaction:strong-acid}{ácido
fraco}.
\end{solution}

\begin{solution}{pb:b1:predominant-reaction:1}
\textbf{1.} \ce{H3PO4}/\ce{H2PO4-}, \ce{H2PO4-}/\ce{HPO4^{2-}},
\ce{HPO4^{2-}}/\ce{PO4^{3-}}; por exemplo, \ce{H3PO4 <=> H2PO4- + H+}.
\textbf{2.} Fronteiras em 2.15, 7.21 e 12.34.
\textbf{3.} pH 1: \ce{H3PO4}; 5: \ce{H2PO4-}; 9: \ce{HPO4^{2-}}; 13:
\ce{PO4^{3-}}.
\textbf{4.} \ce{H2PO4-} e \ce{HPO4^{2-}}.
\textbf{5.} Cada próton deixa um íon de carga negativa maior, que retém o
próton seguinte com mais força.
\textbf{6.} \ce{H3PO4 + H2O <=> H2PO4- + H3O+}; a fórmula
$\frac12(\mathrm{p}K_a + \mathrm{p}C) = 1.58$ viola $\mathrm{pH} <
\mathrm{p}K_a - 1 = 1.15$: cerca de um quinto do ácido reage.
\textbf{7.} $x^2 + K_{a1}x - K_{a1}C = 0$ com $K_{a1} = 10^{-2.15}$: $x =
\qty{0.0233}{mol/L}$, $\mathrm{pH} = 1.63$.
\textbf{8.} \ce{2H2PO4- <=> H3PO4 + HPO4^{2-}}, $K = 10^{2.15 - 7.21} =
10^{-5.06}$.
\textbf{9.} $\mathrm{pH} = \frac12(2.15 + 7.21) = 4.68$.
\textbf{10.} $\mathrm{pH} = \frac12(7.21 + 12.34) = 9.78$.
\textbf{11.} Nenhuma: 1.63, 4.68 e 9.78; é preciso uma mistura de duas
delas.
\textbf{12.} \ce{H3PO4 + OH- -> H2PO4- + H2O}, $K = 10^{14.00 - 2.15} =
10^{11.85}$: \omterm{def:b1:extent-q-and-k:quantitative}{quantitativa}.
\textbf{13.} \qty{0.100}{mol/L} de \ce{H2PO4-}: pH 4.68.
\textbf{14.} Os \qty{0.050}{mol} seguintes de hidróxido convertem a
metade dele: \qty{0.050}{mol/L} de \ce{H2PO4-} e de \ce{HPO4^{2-}}, pH 7.21.
\textbf{15.} \qty{0.100}{mol/L} de \ce{HPO4^{2-}}: pH 9.78.
\textbf{16.} O pH sobe a partir de 1.63, dá um salto em torno de
$n = 0.100$ (até 4.68), sobe lentamente passando por 7.21 em
$n = 0.150$, dá um salto em torno de $n = 0.200$ (até
9.78), passa por 12.34 perto de $n = 0.250$ e chega a cerca de 12.6 em $n = 0.300$ (\qty{0.100}{mol/L} de \ce{PO4^3-}: $x^2/(0.100 - x) = 10^{-1.66}$, $x = 0.037$, pH 12.57).
\textbf{17.} Ela contém quantidades iguais do ácido e da base de um par de
$\mathrm{p}K_a$ 7.21 (\cref{prop:b1:predominant-reaction:buffer}).
\textbf{18.} $10^{7.20 - 7.21} = 0.977$.
\textbf{19.} $n(\ce{H2PO4-}) = 0.100 - x$, $n(\ce{HPO4^{2-}}) = x$.
\textbf{20.} $x/(0.100 - x) = 0.977$, $x = \qty{0.0494}{mol}$.
\textbf{21.} O ácido converte \qty{0.0010}{mol} de \ce{HPO4^{2-}}:
$\mathrm{pH} = 7.21 + \log(0.0484/0.0516) = 7.18$, uma variação de
$-0.02$. Em água pura: de 7.00 para 3.00.
\textbf{22.} Não: o pH depende da razão entre as quantidades, que a
diluição não altera (enquanto as \omterm{def:b1:extent-q-and-k:activity}{atividades} permanecerem próximas das
concentrações).
\textbf{23.} $0.100 + 0.0494 = \qty{0.149}{mol}$ de hidróxido: \textbf{\qty{149}{mL}}
da solução a \qty{1.00}{mol/L}, completados até \qty{1.00}{L}.
\end{solution}
